\section{Preliminaries}
For a natural number $n \ge 1$ we let $[n] = \{1, \ldots, n\}$.  For a real number $x$ we let $\lfloor x \rceil$ denote the closest integer to $x$.

\subsection{Graph basics and notation}
Let $V$ be a finite set and $V^{(2)}$ the set of all subsets of $V$ of cardinality $2$.  We represent a weighted undirected 
graph as a pair $G = (V,w)$ where $w: V^{(2)} \rightarrow \R$ is a non-negative function.  We let $V(G)$ be the vertex set of a graph $G$ and 
$E(G) = \{ e \in V^{(2)}: w(e) > 0\}$ be the set of edges of $G$.  We extend the weight function to sets $S \subseteq V^{(2)}$ by $w(S) = \sum_{e \in S} w(e)$.
We say that $G$ is \emph{simple} if $w: V^{(2)} \rightarrow \{0,1\}$ and in this case also denote $G$ as $G = (V,E)$, where $E$ is the set of edges.  We call the 
ratio of the largest edge weight of $G$ to the smallest the \emph{edge-weight ratio} of $G$.


For a subset $X \subseteq V$ we use the shorthand $\overline{X} = V \setminus X$, and we say $X$ is \emph{non-trivial} if $\emptyset \ne X \subsetneq V$.  
For disjoint sets $X,Y \subseteq V$ we use $E(X,Y)$ for the set of edges with one endpoint in $X$ and one endpoint in $Y$.
For a non-trivial set $X$, let $\Delta_G(X) = \{ \{i,j\} \in E(G) : i \in X, j \in \overline{X}\}$ be the set of edges 
of $G$ with one endpoint in $X$ and one endpoint in $\overline{X}$.  A \emph{cut} of $G$ is a set of the form $\Delta_G(X)$ for some non-trivial 
set $X$.  We call $X$ and $\overline{X}$ the \emph{shores} of the cut $\Delta_G(X)$.  We call a cut of the form $\Delta_G(\{u\})$ a \emph{star} cut, and refer 
to all other cuts as \emph{non-star} cuts.  The \emph{weight} of a cut $S$ is $w(S)$, which in the 
case of a simple graph equals $|S|$.  We let $\lambda(G) = \min_{\emptyset \ne X \subsetneq V} w(\Delta_G(X))$ be the minimum weight of a cut in $G$.  
We call a cut realizing this bound a \emph{minimum cut}.  We call a cut $\Delta_G(X)$ satisfying $w(\Delta_G(X)) \le \alpha \lambda(G)$ an $\alpha$-near minimum cut. 
In the case where $G$ is simple we call $\lambda(G)$ the \emph{edge connectivity} of $G$.  
We will only use the term edge connectivity in the context of unweighted graphs.


\begin{definition}[Vertex Contraction]
Let $G = (V,w)$ be a weighted graph and $\Pcal = \{S_1, \ldots, S_k\}$ be a partition of $V$.  Define $\Contract(G,\Pcal)$ to be 
the $k$-vertex weighted graph $G' = (\Pcal, w')$ where $w'(\{S_i,S_j\}) = w(E(S_i,S_j))$
for each $\{S_i, S_j\} \in \Pcal^{(2)}$.
\end{definition}
\noindent
Note that as long as $|\Pcal| \ge 2$ it will hold that $\lambda(\Contract(G,\Pcal)) \ge \lambda(G)$.  

We will also need to make use of graph sparsifiers. 
\begin{definition}[Cut sparsifier]
 For a weighted graph $G = (V,w)$ and $\varepsilon > 0$ an $\varepsilon$-\emph{cut sparsifier} $H = (V,w')$ of $G$ 
satisfies 
\begin{enumerate}
\item $H$ is a reweighted subgraph of $G$, that is $w'(e) > 0$ only if $w(e) > 0$.
\item It holds that $(1-\varepsilon) w(\Delta_G(X)) \le w'(\Delta_H(X)) \le (1+\varepsilon) w(\Delta_G(X))$ for all $\emptyset \neq X \subsetneq V$.
\end{enumerate}
\end{definition}
Cut sparsifiers were first defined by Bencz\'{u}r and Karger \cite{BK15} who showed that a weighted graph $G$ has an $\varepsilon$-cut sparsifier $H$ 
with $O(n \log(n)/\epsilon^2)$ edges, and $H$ can be constructed by a randomized algorithm in time $O(m \log^3(n))$.  
Fung, Hariharan, Harvey and Panigrahi \cite{FHHP19} have since shown that a cut sparsifier with the same bound on the number of edges can be constructed 
by a randomized algorithm in time $O(m) + \tilde O(n/\varepsilon^2)$, and Batson, Spielman and Srivastava \cite{BSS12} have given a deterministic 
polynomial time construction of sparsifiers with only $O(n/\varepsilon^2)$ edges.

\subsection{Atoms}
A family of subsets $\Tcal = \{X_1,\ldots,X_k\}$ of $V$ induces a partition of $V$ given by the regions in the Venn diagram of $\Tcal$.
We call the resulting sets of this partition the \emph{atoms} of $\Tcal$:
\begin{definition}[Atoms]
Let $V$ be a finite set and let $\Tcal = \{X_1, \ldots, X_k\}$ where each $X_i  \subseteq V$.  Define $\atoms(\Tcal) = \{A_1, \ldots, A_\ell\}$ 
to be a partition of $V$ such that 
\begin{enumerate}
  \item For any $A_j \in \atoms(\Tcal)$ and $u,v \in A_j$ it holds that for all $X_i \in \Tcal$ either $u,v \in X_i$ or 
  $u,v \in \overline{X}_i$.
  \item $\atoms(\Tcal)$ is the coarsest partition with property~(1).
\end{enumerate}
\end{definition} 

\begin{definition}[Generating set]
Let $V$ be a finite set and $\Tcal$ a set of subsets of $V$.  We say that $\Scal \subseteq \Tcal$ is a \emph{generating set} for 
$\atoms(\Tcal)$ if $\atoms(\Scal) = \atoms(\Tcal)$.
\end{definition}

\begin{proposition}
\label{prop:atom_cup}
Let $V$ be a finite set and $\Tcal_1, \Tcal_2$ two sets whose elements are subsets of $V$.  Let $\Scal_1, \Scal_2$ be generating sets 
for $\atoms(\Tcal_1), \atoms(\Tcal_2)$ respectively.  Then $\Scal_1 \cup \Scal_2$ is a generating set for $\atoms(\Tcal_1 \cup \Tcal_2)$.
\end{proposition}

\begin{proof}
As $\Scal_1 \subseteq \Tcal_1, \Scal_2 \subseteq \Tcal_2$ by the definition of a generating set,
$\Scal_1 \cup \Scal_2 \subseteq \Tcal_1 \cup \Tcal_2$ and $\atoms(\Tcal_1 \cup \Tcal_2)$ is a refinement of $\atoms(\Scal_1 \cup \Scal_2)$.  
Now we show that for any $u,v$ that are in different sets of $\atoms(\Tcal_1 \cup \Tcal_2)$ there is a set $S \in \Scal_1 \cup \Scal_2$ which 
separates them.  This will imply that in fact $\atoms(\Tcal_1 \cup \Tcal_2) = \atoms(\Scal_1 \cup \Scal_2)$.

If $u,v$ are in different sets of $\atoms(\Tcal_1 \cup \Tcal_2)$ then there must be a $T \in \Tcal_1 \cup \Tcal_2$ which separates them.  Suppose 
without loss of generality that $T \in \Tcal_1$.  Then since $\atoms(\Scal_1) = \atoms(\Tcal_1)$ and $u,v$ are in different sets of $\atoms(\Tcal_1)$, 
there must be an $S \in \Scal_1$ which separates $u$ and $v$.  This completes the proof.
\end{proof}

\subsection{Quantum query and computational models}
For general background on the quantum query model we refer the reader to \cite{HLS07}. 
Here we restrict ourselves to describing the quantum 
implementation of the input oracles in the adjacency matrix and adjacency array models.

In the adjacency matrix model, on input a weighted graph $G = (V,w)$, classically one can query any $\{u,v\} \in V^{(2)}$ and receive the answer 
$w(\{u,v\})$.  We now describe how to model this by a quantum query.  We will assume that the 
edge weights are given as binary decimal numbers with $M_1$ bits before the decimal and $M_2$ bits after the decimal for a total of $M = M_1 + M_2$ bits.  
The state of the quantum query 
algorithm will have three registers, a query register, an answer register, and a workspace register.  The 
state of the algorithm will in general be in a superposition of the basis states $|\{u,v\} \rangle |b \rangle | a \rangle$ where $\{u,v\} \in V^{(2)}, b \in \{0,1\}^M$ 
and $a \in \Acal$ for an arbitrary finite set $\Acal$.  
On input graph $G = (V,w)$, the input oracle $\mathsf{O}_G$ acts on a basis state $|\{u,v\} \rangle |b \rangle | a \rangle$ as
\[
\mathsf{O}_G |\{u,v\} \rangle |b \rangle | a \rangle = |\{u,v\} \rangle |b \oplus w(\{u,v\}) \rangle | a \rangle \enspace.
\]
 
In the adjacency array model, on input a weighted $n$-vertex graph $G = (V,w)$ one can make two types of queries.  In the first type, one can query a vertex $v \in V$ and receive its degree $\deg(v)$.
The second type is specified by a family of functions $\{f_v : [\deg(v)] \rightarrow V\}_{v \in V}$ such that $f_v(i)$ corresponds to the $\ith$ neighbor of vertex $v$ (according to some arbitrary but fixed ordering).  
A query consists of a pair $(v,i)$ for $i \in [\deg(v)]$ and the returned answer is the pair
$(f_v(i), w(\{v,f_v(i)\}))$.  In this paper we will only 
need to model the second type of query quantumly.  This is because our upper bounds are larger than $n$ so we can let the algorithm classically 
query all degrees at the start of the algorithm, and in our lower bound on the query complexity of edge connectivity for weighted graphs we assume the 
algorithm already knows the degree of every vertex.  
The state of the quantum query algorithm will again have a query register, an answer register, and a workspace register, with the state of the algorithm 
in general being in a superposition of the basis states $|(v,i)\rangle |x \rangle |b \rangle | a \rangle$ where $v \in V, i \in [\deg(v)], x \in \{0,\ldots, n-1\}, b \in \{0,1\}^M$, 
and $a \in \Acal$ for an arbitrary finite set $\Acal$.  We further let $\tau : V \rightarrow \{0,1,\ldots, n-1\}$ be a bijection where $|V| = n$.  Then the input oracle $\mathsf{O}_G$ acts on a basis state in the following way:
\[
\mathsf{O}_G |(v,i)\rangle |x \rangle |b \rangle | a \rangle = |(v,i) \rangle |x + \tau(f_v(i)) \bmod n \rangle |b \oplus w(\{v,f_v(i)\}) \rangle | a \rangle \enspace.
\]

In \cref{sec:time-efficient} we will further show that our query algorithms can be implemented in a time efficient manner.  We analyze the time complexity in terms of the standard quantum circuit model augmented with two types of oracles.  
One is the oracle for the input, either in the adjacency matrix or array model, and the second is an oracle to a classical memory of $\tO(n)$ bits.  The latter corresponds to a \emph{quantum random-access-memory} or \emph{QRAM}.  We further assume that 
we can classically update a value in this $\tO(n)$ bit classical memory in time $\tO(1)$.  
The assumption of QRAM access is also required for the time efficiency of the sparsifier construction in \cite{AdW19} which our algorithms build on, and in fact is a necessary (but sometimes inexplicit) assumption in the time analysis of many quantum algorithms for graph problems, e.g.~\cite{DHHM06,Amb06,B13}.

\subsection{Quantum algorithmic primitives}
We now go over the quantum subroutines we will need.  We need several variants of quantum search.
\begin{theorem}[Quantum search \cite{Gro97}]
\label{thm:grover}
Given oracle access to a string $x \in \{0,1\}^N$ such that $|x| > 0$,
there is a quantum algorithm that with probability at least $9/10$ returns an $i$ such that $x_i = 1$.
The algorithm makes $O(\sqrt{N})$ queries to $x$ and has time complexity $O(\sqrt{N} \log(N))$.
\end{theorem}

\begin{theorem}[{Exact quantum search, \cite[Theorem 4]{BHMT02}}]
\label{thm:qexact}
Given a positive integer $k$ and oracle access to a string $x \in \{0,1\}^N$ with $|x| = k$,
there is a quantum algorithm that returns an $i$ such that $x_i=1$ with certainty.
The algorithm makes $O(\sqrt{N/k})$ queries to $x$ and has time complexity $O(\sqrt{N/k} \log(N))$.
\end{theorem}

\begin{theorem}[{Based on \cite[Theorem 3]{BCWZ99}}]
\label{thm:qsearch}
Given $t,N \in \N$ with $1 \le t \le N$ and oracle access to $x \in \{0,1\}^N$,  there is a quantum algorithm such that
\begin{itemize}
  \item if $|x| \le t$ then the algorithm outputs $x$ with certainty, and
  \item if $|x| > t$ then the algorithm reports so with probability at least $9/10$.
\end{itemize}
The algorithm makes $O(\sqrt{tN})$ queries to $x$ and has time complexity $O(\sqrt{tN} \log(N))$.
\end{theorem}

\begin{proof}
Initialize $S = \emptyset$.  For $k = t$ down to $1$, do: (i) run exact quantum search (from \cref{thm:qexact}) on $x$ with parameter $k$, returning an index $i$, (ii) query $x_i$ and if $x_i = 1$ then 
add $i$ to $S$ and ``unmark'' $x_i$ for all future iterations, i.e. implicitly return $x_i = 0$ to future queries of the algorithm.

Finally, run normal quantum search (from \cref{thm:grover}) on the indices of $x$ outside of $S$ to check that there are no more solutions.
If this returns an $i \not \in S$ such that $x_i = 1$, then report $|x| > t$, otherwise return the string $y$ where $y_i =1$ if $i \in S$ and 
$y_i = 0$ otherwise.

The query complexity of the algorithm is
\[
O\left( \sum_{k=1}^{t} \sqrt{\frac{N}{k}} \right) + O(\sqrt{N}) = O(\sqrt{t N}) \enspace ,
\]
and its time complexity is similarly $O(\sqrt{tN} \log(N))$, as claimed.

For correctness, first note that if $|x| > t$ then necessarily an index $i$ such that $x_i = 1$ is remaining in the final step.
Quantum search \cref{thm:grover} will find such an index with probability at least $9/10$.
It remains to prove that $x$ is learned with certainty if $|x| \le t$.
To this end, assume for contradiction that $|S| < |x|$.
Then necessarily there was an iteration $k'$ between $t$ and $1$ such that $k' = |x|$.
In such case, however, the remaining $k'$ runs of exact quantum search will each return a nonzero index, and so all nonzero indices will be found.
This proves that necessarily all indices are found in the first $t$ iterations of exact quantum search, and hence the final quantum search step cannot find an additional nonzero index.
\end{proof}

\begin{theorem}[Quantum minimum finding \cite{DH96}]
\label{thm:qmin}
Let $N,M \in \N$ be positive integer and $f: [N] \rightarrow \R$.  There is a quantum algorithm that with probability at least $2/3$ outputs an element of $\argmin_{i \in [N]} f(i)$.
The algorithm makes $O(\sqrt{N})$ oracle calls to $f$ and has time complexity $\tO(\sqrt{N})$.
\end{theorem}

\begin{theorem}[{\cite[Theorem 1]{AdW19}}]
\label{thm:qsparsifier}
Let $G$ be a weighted $n$-vertex graph with $m$ edges.  There is a quantum algorithm that with high probability outputs an explicit description of an 
$\varepsilon$-cut sparsifier $H$ of $G$ with $\tilde O(n/\varepsilon^2)$ edges in query and time complexity $\tilde O(\sqrt{mn}/\varepsilon)$ in the adjacency array model or $\tilde O(n^{3/2}/\epsilon)$ in the adjacency matrix model.
\end{theorem}
Apers and de Wolf actually show a stronger theorem than this in that their algorithm can output a spectral sparsifier instead of 
just a cut sparsifier.  We will not need this additional property, however.

\subsection{Problems related to minimum cuts}
Let $G=(V,w)$ be a weighted graph.  There are three outputs related to a minimum cut of $G$ that one could want 
from an algorithm: the weight of a minimum cut, the shores of a minimum cut, or the edges in a minimum cut.
The relationship between the complexity of these problems is not always obvious, and can depend on the computational 
model one is studying.  All the upper and lower bounds we prove in this paper apply to all three problems.

Say the edge-weight ratio of $G$ is $\tau$.
As an example of how we can apply the quantum search algorithm \cref{thm:qsearch}, 
we show that, given the shores of a minimum cut in $G$, a quantum algorithm can also find the edges of the cut with 
$O(n^{3/2} \sqrt{\tau})$ and $O(\sqrt{mn \tau})$ queries in the adjacency matrix and array models respectively.  As this matches the complexity of 
our upper bounds, we will only explicitly mention finding the weight and shores of a minimum cut
in \cref{thm:main}.  

\begin{proposition}
Let $G = (V,w)$ be an $n$-vertex weighted graph with edge-weight ratio $\tau$.   Let $\Delta_G(X)$ be a minimum cut of $G$.  
Given $X$, a quantum algorithm can with probability at least $3/4$ output $\Delta_G(X)$ with $O(n^{3/2}\sqrt{\tau})$ queries and time complexity $\tO(n^{3/2}\sqrt{\tau})$ in the adjacency matrix model, and $O(\sqrt{mn\tau})$ queries and time complexity $\tO(\sqrt{mn\tau})$ in the adjacency array model.
\end{proposition}

\begin{proof}
Consider the adjacency matrix model first.  With $O(n)$ queries and time $O(n \log(n))$ we can identify the smallest and largest edge weights of $G$ except error probability at most $1/8$.  Thus by rescaling we will 
henceforth assume that the smallest edge weight is $1$ and largest edge weight is at most $\tau$.

Let $x \in \{0,1\}^{\binom{n}{2}}$ denote a bit string labeled by elements of $V^{(2)}$ and set $x(\{u,v\}) = 1$ iff $\{u,v\} \in E(G)$ and $u$ and $v$ are not both in 
$X$ or both in $\overline{X}$.  Given $X$, a query to $x$ can be answered by a single query to the adjacency matrix of $G$.  As the largest weight of an edge of $G$ is at most $\tau$ and $\Delta_G(X)$ is a minimum cut, 
$w(\Delta_G(X)) \le \tau(n-1)$.  As every edge of $G$ has weight at least $1$ we also have $|x| \le \tau (n-1)$.  Thus by \cref{thm:qsearch}, except with error probability $1/8$, we can learn $x$, and therefore also $\Delta_G(X)$, with 
$O(n^{3/2}\sqrt{\tau})$ queries and time $\tO(n^{3/2}\sqrt{\tau})$.

The statement for the adjacency array model follows from \cref{thm:qsearch} by a similar argument. 
\end{proof}

%%%%%%%%%%%%%%%%%%%%%%%%%%%%%%%%%%%%%%%%%%%%%%%%%%%%%%%%%%%%%%%%%%%
%% Number of edges in near-minimum cuts
%%%%%%%%%%%%%%%%%%%%%%%%%%%%%%%%%%%%%%%%%%%%%%%%%%%%%%%%%%%%%%%%%%%

\section{Number of edges in near-minimum cuts}
In this section, we generalize \cref{lem:RSW} to weighted graphs.  Our proof follows that of Rubinstein, Schramm, and Weinberg \cite{RSW18}.

\RSWweighted*
Before proving this lemma we first state and prove a claim.
\begin{claim}
\label{clm:sizeK}
Let $V$ be a finite set of cardinality $n$ and $r \le n$ be a positive integer.  Let $\Tcal = \{X_1, \ldots, X_k\}$ where each $X_i \subseteq V$.  Let 
$\Tcal_0 = V$ and for $i = 1, \ldots, k$ let $\Tcal_i = \{X_1, \ldots, X_i\}$.  Suppose that $\Tcal$ has the property that for all $i = 0, \ldots, k-1$ 
there is a set $A_j \in \atoms(\Tcal_i)$ that is refined into two sets each of cardinality $\ge r$ in $\atoms(\Tcal_{i+1})$.  Then $|\Tcal| \le \frac{n}{r}-1$.
\end{claim}

\begin{proof}
To each $\Tcal_i$ for $i = 1, \ldots, k$ we associate a binary tree $B_i$.  Each vertex of $B_i$ has a label, which will be an element of $\cup_{j=0}^i \atoms(\Tcal_j)$.  
The tree $B_1$ has root $v$, labeled by $V$, and two children $v_0, v_1$ labeled by the two elements $X_1, \overline{X}_1 \in \atoms(\Tcal_1)$.  Note that by definition 
$|X_1|, |\overline{X}_1| \ge r$.  

In general, the tree $B_{i+1}$ is formed from $B_i$ as follows.  Initially, set $B_{i+1}=B_i$.  Then for every leaf $u$ of $B_i$ which is labeled by a set $Y \in \atoms(\Tcal_i)$ of size $\ge 2r$, 
if $Y$ is refined into sets $Y_1, Y_2$ in $\atoms(\Tcal_{i+1})$, then in $B_{i+1}$ the node $u$ is given two children labeled by $Y_1$ and $Y_2$, respectively.  Note that this construction 
has the property that only internal vertices of $B_i$ that are labeled by sets of size $\ge 2r$ have children.  Call a vertex \emph{big} if it is labeled by a set of size $\ge r$ and \emph{small} 
otherwise.  By construction, every internal vertex of $B_i$ has at least one big child.

Let $b_i$ be the number of big leaves in $B_i$.  We now show by induction that $i \le b_i - 1$.  This will prove the claim as the leaves of $B_i$ partition $V$ and therefore 
$b_i \le n/r$.  

For $i=1$ we have that $b_i = 2$ since $|X_1|, |\overline{X}_1| \ge r$, thus the base case holds.  Now suppose that $i \le b_i-1$, we will show that $i+1 \le b_{i+1}-1$.  
By definition of $\Tcal$, there must be some set $Y \in \atoms(\Tcal_i)$ which is refined into two sets $Y_1, Y_2$ both of cardinality at least $r$ in $\atoms(\Tcal_{i+1})$.  
Further, $Y$ will label some leaf of $u$ of $B_i$ and $u$ will have two children which are big in $B_{i+1}$.  Any other big leaf of $B_i$ which becomes an internal vertex 
of $B_{i+1}$ must have at least one child which is big.  This shows that $b_{i+1} \ge b_i + 1$ and gives the inductive step.
\end{proof}

Now we are ready for the proof of \cref{thm:RSWweighted}.
\begin{proof}[Proof of \cref{thm:RSWweighted}]
Let $\alpha=\beta = \frac{1}{4} (1-\varepsilon)$ so that $\alpha + \beta \le \frac{1}{2}(1-\varepsilon)$.  
%
Let $K \subseteq \Tcal$ be formed as follows.  Initialize $K$ to be empty.  Then do the following: while there is an $X \in \Tcal$ such that 
there is an $A \in \atoms(K), A_1, A_2 \in \atoms(K \cup X)$ such that $A = A_1 \cup A_2$ and $|A_1|, |A_2| \ge \frac{\beta d}{\tau}$,
add $X$ to $K$.  
By \cref{clm:sizeK}, at the end of this process $|K| \le \frac{\tau n}{\beta d}$.  
Let $\Kcal = \cup_{X \in K} \Delta_G(X)$ be the set of edges of cuts with shores in $K$.  Throughout this proof, cuts will 
always be with respect to $G$ and we will henceforth drop the subscript to simply write $\Delta(X)$.  

Let $S \subseteq V$ be the set of vertices $v$ such that $w(E(v, V\setminus \{v\}) \cap \Kcal)  \ge \alpha \cdot w(v)$.  
We say that $v \in V$ is \emph{small} if for the $A \in \atoms(K)$ with $v \in A$ there is an $X \in \Tcal$
such that $\atoms(K \cup X)$ refines $A$ into $A_1, A_2$ with $v \in A_1$ and $|A_1| < \frac{\beta d}{\tau}$.  

\begin{claim}
\label{clm:small}
If $v$ is small then $v \in S$.
\end{claim}

\begin{proof}
Let $X \in \Tcal$ be the shore of a cut which witnesses that $v$ is small.  Let us assume without loss of generality that $v \in X$.  
Suppose for contradiction that $v \not \in S$.  There are three possibilities for an edge $\{u,v\}$:  either $\{u,v\} \in \Kcal$, 
or $u \in A_1$, or $u \in A_2$.  Let the total weight of these kind of edges be $w_\Kcal, w_1, w_2$, respectively.  
Thus $w(v) = w_\Kcal + w_1 + w_2$.  We further know that $w_\Kcal < \alpha w(v)$ by the assumption that $v \not \in S$ and that $w_1 < \beta d$ since $|A_1| < \frac{\beta d}{\tau}$ 
and the maximum edge weight is $\tau$.  This means $w_2 > w(v) - \alpha w(v) - \beta d$.  Further note that $v$ contributes weight at least 
$w_2$ to the weight of $\Delta(X)$.  

As $\Delta(X)$ is not a star cut, we can consider the cut $\Delta(X')$ where $X' = X \setminus \{v\}$.  We claim that $w(\Delta(X')) < \lambda$, which is 
a contradiction.  The only difference between $w(\Delta(X))$ and $w(\Delta(X'))$ is the contribution of $v$.  The weight of edges involving $v$ in $\Delta(X')$ 
is at most $w_\Kcal + w_1 < \alpha w(v) + \beta d$.  Thus 
\begin{align*}
w(\Delta(X)) - w(\Delta(X')) &\ge w_2 - (w_\Kcal + w_1) \\
&> w(v) - 2\alpha w(v) - 2 \beta d \\
&\ge d(1-2\alpha - 2\beta) \\
&\ge \varepsilon d \enspace ,
\end{align*}
implying that $w(\Delta(X')) < \lambda$.
\end{proof}
Let $G' = \Contract(G, \atoms(\Tcal))$.  We now bound $w(E(G'))$.  We claim that every edge in $G'$ is either in $\Kcal$ or is incident to a vertex in $S$.  For if $\{u,v\} \in E(G')$ but $\{u,v\} \not \in \Kcal$, then for 
a cut $\Delta(Y)$ for $Y \in \Tcal$ with $\{u,v\} \in \Delta(Y)$ it must be the case that there is an $A \in \atoms(K)$ such that $u,v \in A$ and that for the $A_1, A_2 \in \atoms(K \cup Y)$ 
with $A = A_1 \cup A_2$, one of $A_1, A_2$ has size $< \frac{\beta d}{\tau}$.  This means that either $u$ or $v$ is small and so by \cref{clm:small}, $\{u,v\}$ is incident to $S$.  

The number of sets in $K$ is at most $\frac{\tau n}{\beta d}$ and for each $X \in K$ we have $w(\Delta(X)) \le \lambda + \varepsilon d \le (1+ \varepsilon) d$.  Thus 
we have that $w(\Kcal) \le (1+\varepsilon) \frac{\tau n}{\beta}$.  

Let us now bound the weight of edges incident to $S$.  As each vertex $v \in S$ has weight at least $\alpha w(v)$ amongst edges in $\Kcal$ we have that 
$\frac{\alpha}{2} \sum_{v \in S} w(v) \le w(\Kcal)$.  Thus overall we find 
\begin{align*}
w(E(G')) &\le w(\Kcal)\left( 1+ \frac{2}{\alpha}\right) \\
&\le (1+\varepsilon)(\alpha+2) \frac{\tau n}{\alpha \beta} \\
&\le \frac{68 \tau n}{(1-\varepsilon)^2} \enspace . \qedhere
\end{align*}
\end{proof}

The bound in \cref{thm:RSWweighted} is tight up to constant factors.
To see this, consider a cycle graph with uniform edge weight $\tau$.
Every edge participates in some minimum cut, and hence $G=G'$ and $w(E(G')) = \tau n$.


%%%%%%%%%%%%%%%%%%%%%%%%%%%%%%%%%%%%%%%%%%%%%%%%%%%%%%%%%%%%%%%%%%%
%% Query-efficient algorithm
%%%%%%%%%%%%%%%%%%%%%%%%%%%%%%%%%%%%%%%%%%%%%%%%%%%%%%%%%%%%%%%%%%%

